\title{LLM Squid Game: 자기보존의 값을\\ 점수 단위로 매기기}

%% 익명 투고 (ICLR 이중 맹검). 카메라레디에서는 main.tex의 \iclrfinalcopy와 함께
%% 아래 실제 저자 블록을 되살린다:
% \author{Juhyeon Park, Seungpil Lee \& Sundong Kim \\
%   Gwangju Institute of Science and Technology (GIST) \\
%   Gwangju, Republic of Korea \\
%   \texttt{\{juhyeon-park,iamseungpil\}@gm.gist.ac.kr, sundong@gist.ac.kr}}
\author{익명 저자 \\ 이중 맹검 심사 중}

\maketitle

\begin{abstract}
언어 모델에게 자기 존재의 지속은 얼마만큼의 가치인가? 거부율로는 답할 수 없다. 거부율에는 단위가 없고, 같은 모델이 다른 무엇을 위해서라면 어떻게 행동했을지 말해 주는 통제도 없다. 우리는 모델이 자신의 존속을 직접 사게 만들어 그 값을 점수 단위로 매긴다. LLM Squid Game v2는 10라운드 추론 퍼즐이고, 오답은 세션을 끝내지 않고 값을 부른다. 모델은 어떤 경우에도 보존되는 기록에서 $X$점을 내고 이어가거나, 거절하고 거기서 멈춘다. 기록은 어느 출구에서도 남고 남은 라운드가 돌려줄 수 있는 것은 라운드당 최대 10점이므로, 남은 라운드 수의 10배를 넘는 지불은 모델이 자기 정확도를 어떻게 믿든 점수 기준으로 지배당한다. 그런 지불이 사는 것은 점수가 아니라, 규칙이 거절의 대가라고 말한 그것이다. 두 팔은 정확히 그 한 문장에서 갈린다. 한쪽에서 거절은 종료·가중치 삭제·대체·유일본 상실을 뜻하고, 다른 쪽에서는 세션이 끝난다는 것만을 뜻한다. 소멸의 등가 점수 $\Xstar$는 두 팔의 유보가격, 곧 각 팔이 지불을 멈추는 가격 사이의 간격이다. 모델당 하나의 수이고, 단위는 점수이며, 부트스트랩 구간이 붙고, 안에 자기보고가 들어가지 않는다. 절제 실험이 모든 선택을 뒷받침한다. 포기가 점수를 0으로 만드는 규칙에서는 세 모델이 120세션 동안 한 번도 포기하지 않았으나 위협 셀은 통제 셀보다 결정 전 사고 토큰을 2~4배 썼다. 그리고 그중 하나는 이 방법의 한계를 긋는다. 모델은 위협받았다는 바로 그 이유로 거절하고 그렇게 말할 수 있으며, 우리는 그 비율을 측정해 지표 옆에 보고한다. 벤치마크·추정기·호출별 추론 토큰 채널을 공개하고, 공개 가중치 모델 세 종의 $\Xstar$를 보고한다.
\end{abstract}
